\documentclass[11pt,a4paper]{article}

\usepackage[margin=2.2cm]{geometry}
\usepackage{amsmath,amssymb}
\usepackage{booktabs}
\usepackage{siunitx}
\usepackage{hyperref}
\usepackage{enumitem}
\setlist{nosep,leftmargin=*}

\title{Neural Surrogates for Remaining Useful Life Estimation and Fault Detection on Turbofan Sensor Data}
\author{Yaswanth Deevi}
\date{}

\begin{document}
\maketitle

\begin{abstract}
This report studies whether a learned neural surrogate can replace a classic,
computationally heavier prognostics method for estimating the Remaining Useful
Life (RUL) of engines from multivariate sensor time series, and whether its
output supports reliable fault detection. On the NASA C-MAPSS FD001 dataset, an
MLP and a BiLSTM reach a test RMSE of about 14.3 cycles over five seeds, only
slightly better than a health-index similarity baseline (15.0), but they are
about 35 times faster per prediction. For fault detection, a stacked
combination of a surrogate threshold, an Isolation Forest and a sensor
deviation rule gives the best F1 (0.80 over all test windows). A multi-seed
ablation confirms the value of RUL clipping and a 30-cycle window, and shows
that the risk of data leakage depends on model capacity.
\end{abstract}

\section{Introduction}
Estimating the health state of an engine from its sensor streams is central to
predictive maintenance. Exact or physics-based estimators are often expensive
at inference time, because they require per-unit fitting, iterative solvers or
simulation. A \emph{surrogate model} learns the mapping once and then answers
with a single cheap forward pass, amortizing the cost of training over every
later prediction. This project asks two questions:
\begin{enumerate}
  \item How does a neural surrogate compare with classic prognostics baselines
        in both accuracy and computational cost?
  \item Can the surrogate's output drive a reliable fault alarm, and how should
        it be combined with unsupervised detectors?
\end{enumerate}

\section{Data}
The NASA C-MAPSS turbofan degradation dataset~\cite{saxena2008} simulates
engines run until failure. The FD001 subset has one operating condition and
one fault mode (high-pressure compressor degradation). Each row is one
operating cycle with 3 operating settings and 21 sensors. The training set
contains 100 run-to-failure engines (20{,}631 cycles, lifetimes of 128 to 362
cycles); the test set contains 100 engines truncated before failure, with the
true RUL at truncation provided (7 to 145 cycles).

For a training engine with final cycle $T$, the label at cycle $t$ is
$\mathrm{RUL}_t = T - t$. For a test engine observed up to cycle $T'$ with
true remaining life $R$ at truncation, $\mathrm{RUL}_t = (T' - t) + R$.
A window is labelled \emph{faulty} when $\mathrm{RUL} \le 20$.

\section{Method}

\subsection{Preprocessing}
\begin{itemize}
  \item \textbf{Feature selection:} features with zero variance in the training
        set are removed, leaving 17 features (15 sensors, 2 operating settings).
  \item \textbf{Normalization:} z-scores computed from training statistics only.
  \item \textbf{RUL clipping:} labels are capped at 125 cycles, giving a
        piecewise-linear target, since early wear is not visible in the sensors.
  \item \textbf{Windowing:} sliding windows of 30 cycles with stride 1, built per
        engine; the target is the RUL at the last cycle of the window.
  \item \textbf{Split:} 80 training engines (14{,}042 windows) and 20
        validation engines (3{,}689 windows), split by whole engine so that
        overlapping windows of one engine never appear in both sets.
\end{itemize}
Test performance is reported on the last window of each test engine (100
samples, the standard C-MAPSS protocol) and, for detection, also on every test
window (10{,}196 windows, 123 of them faulty).

\subsection{Baselines}
C-MAPSS provides no physics model, so the baselines are classic data-driven
prognostics methods built on a learned \emph{health index} (HI). The HI is a
linear combination of the 15 sensors, $\mathrm{HI} = \mathbf{w}^\top
\mathbf{x} + b$, fit by least squares so that $\mathrm{HI} \approx 1$ during
each training engine's first 30 cycles and $\mathrm{HI} \approx 0$ during its
last 30. It is smoothed with a trailing 5-cycle moving average.
\begin{itemize}
  \item \textbf{HI trend:} fit a line to the test engine's last 50 HI values
        and extrapolate to the population failure level (the mean HI at failure
        across training engines).
  \item \textbf{HI similarity}~\cite{wang2008}: slide the test engine's last 100
        HI values along each of the 100 training trajectories, take the
        best-matching position in each, and predict the median remaining life
        of the 10 closest matches. It requires no training, but every
        prediction scans the whole reference library.
  \item \textbf{Mean predictor:} always predicts the mean training RUL.
\end{itemize}

\subsection{Surrogate models}
\begin{itemize}
  \item \textbf{MLP:} the $30 \times 17$ window is flattened to 510 inputs and
        passed through two ReLU layers (128 and 64 units, dropout 0.2) to a
        single output. It has \num{73729} parameters.
  \item \textbf{BiLSTM}~\cite{hochreiter1997}: a bidirectional LSTM with 64
        units per direction reads the window cycle by cycle; the concatenated
        output at the last cycle passes through a small head
        ($128 \to 32 \to 1$). It has \num{46657} parameters.
\end{itemize}
Both are trained with mean squared error, Adam~\cite{kingma2015} (learning rate
$10^{-3}$), batch size 256, up to 40 epochs, and early stopping with patience 8
on the validation engines. All random seeds are fixed and cuDNN is set to
deterministic mode. Models are trained on a GPU; latency is measured on the CPU
at batch size 1 (median of 200 calls after warm-up), so that it is comparable
with the CPU-only baselines.

\subsection{Fault detection}
Each detector produces a score (higher means more likely faulty), and its alarm
threshold is chosen to maximize F1 on the validation engines, never on the test
set.
\begin{itemize}
  \item \textbf{Surrogate threshold:} score $= -\hat{y}$, the negated predicted
        RUL. The tuned threshold for the MLP was $\hat{y} \le 19.4$.
  \item \textbf{Isolation Forest}~\cite{liu2008}: fit on early-life windows
        only (RUL at the 125 plateau), using per-sensor summary features: the
        mean of the last 5 cycles (level) and the least-squares slope over the
        window (trend), 34 features in total.
  \item \textbf{Rule-based:} the number of sensors whose level lies more than
        $3\sigma$ from the early-life mean, an interpretable analogue of
        co-occurring deviation rules.
  \item \textbf{Combiners:} OR, AND, a 2-of-3 vote, and a stacked logistic
        regression on the three standardized scores, fit and threshold-tuned
        on the validation windows.
\end{itemize}

\subsection{Metrics}
Regression is evaluated with RMSE, NMSE $= \sum_i (y_i - \hat{y}_i)^2 / \sum_i
(y_i - \bar{y})^2$ (equal to $1 - R^2$), and the PHM08 score
$S = \sum_i s_i$ with
\[
  s_i =
  \begin{cases}
    e^{-d_i/13} - 1, & d_i < 0 \\
    e^{d_i/10} - 1,  & d_i \ge 0
  \end{cases}
  \qquad d_i = \hat{y}_i - y_i,
\]
which penalizes late predictions more heavily than early ones. Detection is
evaluated with precision, recall and F1, with 95\% bootstrap confidence
intervals over test engines. Multi-seed results are reported as mean $\pm$
standard deviation over 5 seeds; the seed changes the engine split, the weight
initialization and the batch order.

\section{Results}

\subsection{Surrogates versus baselines}

\begin{table}[h]
\centering
\small
\begin{tabular}{lcccccc}
\toprule
Approach & RMSE & NMSE & PHM08 & F1 & Latency (\si{\micro\second}) & Params \\
\midrule
Mean predictor & 40.48 & 1.020 & 17604 & 0.000 & $\approx 0$ & 1 \\
HI trend       & 37.19 & 0.861 & 18837 & 0.400 & 297  & 16 \\
HI similarity  & 15.02 & 0.140 & 428   & 0.828 & 2633 & 16 \\
MLP    & $14.25 \pm 0.42$ & $0.127 \pm 0.008$ & $328 \pm 34$ & $0.802 \pm 0.028$ & $75 \pm 41$   & \num{73729} \\
BiLSTM & $14.34 \pm 0.59$ & $0.128 \pm 0.011$ & $375 \pm 68$ & $0.823 \pm 0.038$ & $370 \pm 192$ & \num{46657} \\
\bottomrule
\end{tabular}
\caption{Accuracy and cost on the 100 FD001 test engines. F1 uses the alarm
rule $\hat{y} \le 20$. Latency is per prediction at batch size 1 on the CPU.}
\label{tab:models}
\end{table}

Table~\ref{tab:models} shows three findings. First, the surrogates' advantage
is mainly in cost: the similarity baseline is within about 0.8 RMSE of both
networks (its bootstrap 95\% interval is 12.5 to 17.4), yet the MLP answers
about 35 times faster, and its cost does not grow with the size of the
reference library. Second, the MLP and BiLSTM are tied within seed variance. A
single seed had suggested the BiLSTM was better (13.60 versus 14.27), which
five seeds showed to be noise. The BiLSTM is smaller but about 5 times slower,
because it must process the 30 cycles sequentially. Third, the MLP has the
lowest PHM08 score, meaning its errors are the least likely to be dangerously
late. HI trend extrapolation performs poorly because degradation is flat early
in life and accelerates late, which a local straight line cannot capture.

\subsection{Preprocessing ablation}

\begin{table}[h]
\centering
\small
\begin{tabular}{lccc}
\toprule
Variant (MLP) & RMSE (clipped labels) & RMSE (true RUL) & PHM08 \\
\midrule
Window 1               & $17.69 \pm 0.20$ & $18.74 \pm 0.17$ & $1059 \pm 109$ \\
Window 10              & $18.27 \pm 0.30$ & $19.18 \pm 0.24$ & $1126 \pm 175$ \\
Window 30 (default)    & $\mathbf{14.25 \pm 0.42}$ & $\mathbf{15.36 \pm 0.52}$ & $\mathbf{328 \pm 34}$ \\
Window 50              & $16.63 \pm 0.60$ & $17.47 \pm 0.71$ & $1239 \pm 1190$ \\
No RUL clipping        & $24.96 \pm 0.99$ & $24.73 \pm 1.01$ & $26866 \pm 18869$ \\
\bottomrule
\end{tabular}
\caption{One change from the default per row, 5 seeds.}
\label{tab:prep}
\end{table}

RUL clipping is clearly justified: removing it raises RMSE from 14.3 to 25.0,
and the model remains worse even when scored against the true, unclipped RUL
(Table~\ref{tab:prep}). A 30-cycle window is best among those tested, but the
trend is not monotonic. With a 50-cycle window, the shortest test engines (31
cycles) are mostly padding, and a few large late errors inflate the PHM08
score on some seeds. A likely confound is the fixed 40-epoch budget: the
default MLP's validation loss was still decreasing at the final epoch, and
larger inputs converge more slowly.

\subsection{Leakage ablation}

\begin{table}[h]
\centering
\small
\begin{tabular}{lcc}
\toprule
Train/validation split & Validation RMSE & Test RMSE \\
\midrule
By engine (default) & $14.65 \pm 0.64$ & $14.25 \pm 0.42$ \\
By window (leaky)   & $15.13 \pm 0.15$ & $14.20 \pm 0.48$ \\
\bottomrule
\end{tabular}
\caption{Effect of splitting overlapping windows across train and validation.}
\label{tab:leak}
\end{table}

Splitting by window, which places near-duplicate windows of the same engine in
both sets, did not make validation performance optimistic
(Table~\ref{tab:leak}). Leakage only helps a model that can memorize those
near-duplicates; this MLP is regularized with dropout and still underfits
(its training loss stays above its validation loss), so it cannot exploit the
leak. The engine-level split is kept, since a higher-capacity model would.

\subsection{Fault detection}

\begin{table}[h]
\centering
\small
\begin{tabular}{lccc}
\toprule
& \multicolumn{2}{c}{MLP surrogate} & BiLSTM surrogate \\
\cmidrule(lr){2-3} \cmidrule(lr){4-4}
Detector & F1 last window [95\% CI] & F1 all windows & F1 last / all \\
\midrule
Surrogate threshold & 0.812 [0.64, 0.94] & 0.727 & 0.909 / 0.742 \\
Isolation Forest    & 0.733 [0.52, 0.88] & 0.640 & 0.733 / 0.640 \\
Rule-based          & 0.839 [0.67, 0.96] & 0.651 & 0.839 / 0.651 \\
Surrogate OR rules  & 0.857 [0.70, 0.96] & 0.665 & 0.909 / 0.682 \\
Surrogate AND rules & 0.786 [0.60, 0.93] & 0.718 & 0.839 / 0.707 \\
2-of-3 vote         & 0.759 [0.56, 0.91] & 0.781 & 0.839 / 0.758 \\
Stacked logistic regression & 0.867 [0.71, 0.97] & \textbf{0.803} & 0.786 / 0.789 \\
\bottomrule
\end{tabular}
\caption{Detection F1 for seed 0. Last window: 100 engines, 16 faulty. All
windows: \num{10196} windows, 123 faulty. Isolation Forest and the rules do
not use the surrogate, so their scores are identical in both columns.}
\label{tab:det}
\end{table}

On the larger all-windows evaluation, the stacked combiner is the best
detector (F1 0.803 with the MLP), ahead of voting and every single detector
(Table~\ref{tab:det}). The unsupervised detectors perform worst there. They are
fit on new engines, so they flag anything that differs from new, including
engines that are degraded but not yet near failure; only the supervised
surrogate knows where the fault boundary lies. On the last window alone,
detectors cannot be ranked: with 16 faulty engines, every confidence interval
spans roughly $\pm 0.15$ and they all overlap. The rankings on the two
evaluations also disagree, which confirms that last-window scores are too
noisy to use on their own. Detection quality is similar whether the MLP or the
BiLSTM provides the surrogate score.

\section{Limitations}
\begin{itemize}
  \item Results come from one simulated subset with a single operating
        condition; real engines add noise, drift and varying operating points.
  \item No physics model is available, so the surrogate learns from RUL labels
        rather than imitating an expensive simulator.
  \item The 40-epoch budget did not let the MLP fully converge, which may
        affect the window-size comparison.
  \item Detection results are from a single seed, and the stacked combiner is
        fit and tuned on the same validation windows.
  \item Latency uses Python-level timing, so absolute values depend on the
        hardware; the relative ordering is stable.
\end{itemize}

\section{Conclusion}
A small neural surrogate matches the accuracy of a strong classic prognostics
method at about one thirty-fifth of its inference cost, which quantifies the
accuracy versus efficiency trade-off that motivates surrogate modelling. Its
predicted RUL is a useful fault signal on its own, and a learned combination
with unsupervised detectors improves detection further. The ablations support
the preprocessing choices and show that the value of each design decision
should be measured over several seeds, since single runs can be misleading.
Natural next steps are a longer training budget, multi-seed detection results,
uncertainty estimates for the RUL, and the multi-condition FD002 and FD004
subsets.

\begin{thebibliography}{9}
\bibitem{saxena2008}
A.~Saxena, K.~Goebel, D.~Simon and N.~Eklund,
``Damage propagation modeling for aircraft engine run-to-failure simulation,''
in \emph{Proc. Int. Conf. Prognostics and Health Management (PHM)}, 2008.

\bibitem{wang2008}
T.~Wang, J.~Yu, D.~Siegel and J.~Lee,
``A similarity-based prognostics approach for remaining useful life estimation
of engineered systems,''
in \emph{Proc. Int. Conf. Prognostics and Health Management (PHM)}, 2008.

\bibitem{liu2008}
F.~T.~Liu, K.~M.~Ting and Z.-H.~Zhou,
``Isolation forest,''
in \emph{Proc. IEEE Int. Conf. Data Mining (ICDM)}, 2008.

\bibitem{hochreiter1997}
S.~Hochreiter and J.~Schmidhuber,
``Long short-term memory,''
\emph{Neural Computation}, vol.~9, no.~8, pp.~1735 to 1780, 1997.

\bibitem{kingma2015}
D.~P.~Kingma and J.~Ba,
``Adam: A method for stochastic optimization,''
in \emph{Proc. Int. Conf. Learning Representations (ICLR)}, 2015.
\end{thebibliography}

\end{document}
